% !TEX root = task-2-2-3.tex
% Editable task bank. Append new BankTask blocks at the end of the document.
\documentclass[12pt,a4paper]{article}
\usepackage{physics-textbook}
\usepackage{graphicx,needspace}
\usetikzlibrary{decorations.pathmorphing}
\newif\ifShowSolutions
\ifdefined\TutorEdition
  \ShowSolutionstrue
\else
  \ShowSolutionsfalse % Change to \ShowSolutionstrue to display solutions.
\fi
\newcounter{banktask}
\newsavebox{\BankTaskBox}
\newcommand{\BankTaskContent}[3]{%
  \noindent\textbf{\thebanktask.}\quad #2\par
  \ifShowSolutions
    \Needspace{50mm}% Keep the solution heading with its opening text or diagram.
    #3\par
    \ifdefined\TutorEdition
      \ifcsname TutorSource#1\endcsname
        \nopagebreak\smallskip{\small\csname TutorSource#1\endcsname\par}%
      \fi
    \fi
  \fi
}
\newcommand{\BankTask}[3]{%
  \refstepcounter{banktask}\par\medskip
  \sbox{\BankTaskBox}{\begin{minipage}{\linewidth}%
    \BankTaskContent{#1}{#2}{#3}\end{minipage}}%
  \ifdim\dimexpr\ht\BankTaskBox+\dp\BankTaskBox\relax>\dimexpr\textheight-20mm\relax
    \Needspace{35mm}\BankTaskContent{#1}{#2}{#3}%
  \else
    \Needspace{\dimexpr\ht\BankTaskBox+\dp\BankTaskBox\relax}%
    \noindent\usebox{\BankTaskBox}\par
  \fi
  \medskip
}
\begin{document}
\BookSubsection{2.2.3}{Жидкости}


% Source sdamgia-828: https://phys-ege.sdamgia.ru/problem?id=828
\BankTask{sdamgia-828}{
В процессе перехода вещества из жидкого состояния в кристаллическое

$1$) существенно увеличивается расстояние между его молекулами

$2$) молекулы начинают притягиваться друг к другу

$3$) существенно увеличивается упорядоченность в расположении его молекул

$4$) существенно уменьшается расстояние между его молекулами
}{
\smallskip\noindent\textbf{Решение}\par
В процессе перехода вещества из жидкого состояния в твердое происходит построение кристаллической решетки, а значит, существенно увеличивается упорядоченность в расположении молекул вещества.

Ответ: $3$.
}

% Source sdamgia-3208: https://phys-ege.sdamgia.ru/problem?id=3208
\BankTask{sdamgia-3208}{
Какие из утверждений справедливы для жидкостей?

А) Характерное расстояние между молекулами жидкости по порядку величины совпадает с размерами самих молекул 

Б) Жидкость сохраняет форму

В) Жидкость сохраняет объем

Г) Молекулы жидкости образуют периодичную решетку

$1$. А и Г 

$2$. Б и Г

$3$. А и В

$4$. Б и В
}{
\smallskip\noindent\textbf{Решение}\par
Молекулы жидкости достаточно плотно упакованы, так что утверждение А верно. В силу плотной упаковки и достаточно сильного взаимодействия между молекулами, жидкости практически несжимаемы. Как следствие, они сохраняют объем. Молекулы жидкости ведут так называемый «кочевой образ жизни», они могут переходить из одного положения в другое, этим обеспечивается текучесть жидкостей, таким образом, жидкости не сохраняют форму. Утверждение Г относится к кристаллическим телам. Окончательно, верными являются утверждения А и В.

Ответ: $3$.
}

% Source sdamgia-3211: https://phys-ege.sdamgia.ru/problem?id=3211
\BankTask{sdamgia-3211}{
Вещество сохраняет объем, но не сохраняет форму. Это утверждение соответствует модели

$1$) только газа

$2$) только жидкости

$3$) только твердого тела

$4$) газа, жидкости и твердого тела
}{
\smallskip\noindent\textbf{Решение}\par
Твердое тело имеет кристаллическую решетку, поэтому оно сохраняет и форму и объем. Газ занимает весь отведенный ему объем, так что не сохраняет ни форму, ни объем. Молекулы жидкости упакованы практически также плотно, как и молекулы кристаллических тел, и взаимодействуют достаточно сильно, поэтому жидкость сохраняет объем. Однако молекулы жидкости не образуют решетки, они ведут так называемый «кочевой образ жизни», что обеспечивает текучесть жидкости. Как результат жидкости не сохраняют форму.

Ответ: $2$.
}

% Source sdamgia-7695: https://phys-ege.sdamgia.ru/problem?id=7695
\BankTask{sdamgia-7695}{
В стакане находится вода в кристаллическом состоянии при температуре $0^\circ$C. Из-за контакта с теплым стаканом и с комнатным воздухом вся вода переходит в жидкое состояние. В процессе перехода воды из кристаллического состояния в жидкое

$1$) ее объем уменьшается;

$2$) ее плотность уменьшается;

$3$) ее внутренняя энергия уменьшается;

$4$) ее температура уменьшается.
}{
\smallskip\noindent\textbf{Решение}\par
При переходе воды из кристаллического состояния в жидкое, ее плотность увеличивается, а значит при неизменной массе вещества объем должен уменьшиться.

Ответ: $1$
}


% Source sdamgia-6301: https://phys-ege.sdamgia.ru/problem?id=6301
\BankTask{sdamgia-6301}{
Плотность воды при температуре $100^\circ$C равна $950\,\text{кг}$/$\text{м}^{3}$, а наибольшая плотность водяного пара при $100^\circ$C равна $0{,}59\,\text{кг}$/$\text{м}^{3}$. Такое различие плотностей связано с тем, что

$1$) молекулы жидкости и пара движутся с разными скоростями

$2$) при переходе молекул из жидкости в пар уменьшается энергия их взаимодействия

$3$) число молекул в $1$ $\text{м}^{3}$ пара меньше чем в $1$ $\text{м}^{3}$ воды

$4$) молекулы жидкости и пара имеют разные массы
}{
\smallskip\noindent\textbf{Решение}\par
Молекулы внутри пара и жидкости находятся на различном расстоянии друг от друга, в жидкостях это расстояние значительно меньше, чем в газах. Поэтому число молекул в $1$ $\text{м}^{3}$ пара меньше чем в $1$ $\text{м}^{3}$ воды.

Ответ: $3$.
}

% Source sdamgia-3287: https://phys-ege.sdamgia.ru/problem?id=3287
\BankTask{sdamgia-3287}{
Явление диффузии в жидкостях объясняется тем, что молекулы жидкостей

$1$) отталкиваются друг от друга

$2$) колеблются около своих положений равновесия

$3$) притягиваются друг к другу

$4$) могут хаотично перемещаться по объему
}{
\smallskip\noindent\textbf{Решение}\par
Диффузией называется процесс взаимного проникновения молекул одного вещества между молекулами другого, приводящий к самопроизвольному выравниванию их концентраций по всему занимаемому объему. Молекулы жидкости могут хаотично перемещаться по объему (они ведут «кочевой образ жизни»). Именно этим и объясняется главным образом явление диффузии в жидкостях.

Ответ: $4$.
}

% Source sdamgia-6193: https://phys-ege.sdamgia.ru/problem?id=6193
\BankTask{sdamgia-6193}{
Если толченый мел размешать в воде, то частицы мела будут долго «висеть» в толще воды, не оседая на дно. Это явление объясняется тем, что

$1$) вода выталкивает их вверх согласно закону Архимеда

$2$) частицы мела совершают броуновское движение в воде

$3$) Земля не притягивает столь мелкие частицы

$4$) температура частиц мела выше температуры воды
}{
\smallskip\noindent\textbf{Решение}\par
Частички толченого мела достаточно маленькие, чтобы действие, оказываемое на них тепловым движением молекул воды было существенным. Следовательно, частицы мела совершают броуновское движение, вследствие чего долго не оседают на дно.

Правильный ответ указан под номером $2$.
}

% Source sdamgia-4485: https://phys-ege.sdamgia.ru/problem?id=4485
\BankTask{sdamgia-4485}{
Диффузия в жидкости происходит быстрее при повышении температуры, потому что с повышением температуры

$1$) увеличиваются силы взаимодействия молекул

$2$) увеличивается скорость теплового движения молекул

$3$) жидкости расширяются

$4$) уменьшаются силы взаимодействия молекул
}{
\smallskip\noindent\textbf{Решение}\par
Диффузия — это процесс взаимного проникновения молекул одного вещества между молекулами другого, приводящий к самопроизвольному выравниванию их концентраций по всему занимаемому объему. Он происходит за счет непрерывного хаотического движения молекул. Как известно, увеличение температуры приводит к увеличению скорости теплового движения. Поэтому правильным является ответ под номером $2$.

Ответ: $2$.
}

% Source sdamgia-3372: https://phys-ege.sdamgia.ru/problem?id=3372
\BankTask{sdamgia-3372}{
На газовой плите стоит высокая кастрюля с водой, закрытая крышкой. Если воду из нее перелить в широкую кастрюлю, у которой площадь дна вдвое больше, и тоже закрыть крышкой, то вода закипит заметно быстрее, чем если бы она осталась в узкой. Этот факт объясняется тем, что

$1$) увеличивается площадь нагревания и, следовательно, увеличивается скорость нагревания воды

$2$) в $2$ раза уменьшается необходимое давление насыщенного пара в пузырьках и, следовательно, воде у дна надо нагреваться до менее высокой температуры

$3$) увеличивается площадь поверхности воды и, следовательно, испарение идет более активно

$4$) в $2$ раза уменьшается глубина слоя воды и, следовательно, пузырьки пара быстрее добираются до поверхности
}{
\smallskip\noindent\textbf{Решение}\par
Процесс кипения начинается при температуре, когда внешнее давление на жидкость сравнивается с давлением насыщенных паров жидкости при этой температуре. От площади поверхности жидкости температура кипения никак не зависит То, что вода закипит в более широкой кастрюле быстрее, объясняется тем, что при этом увеличивается площадь нагревания, и, следовательно, увеличивается скорость нагревания воды.

Ответ: $1$.
}

% Source sdamgia-4343: https://phys-ege.sdamgia.ru/problem?id=4343
\BankTask{sdamgia-4343}{
Модель, служащая для демонстрации внутреннего строения тел, устроена следующим образом: в дно прямоугольной коробки воткнуты одинаковые упругие вертикальные стерженьки, на каждый из которых насажен магнитик в виде плоской таблетки. После приведения одного из магнитиков в колебательное движение вскоре начинают хаотически колебаться на стерженьках и все остальные магнитики, отталкиваясь друг от друга. Эта модель лучше всего иллюстрирует поведение молекул

$1$) идеального газа

$2$) жидкости

$3$) твердого тела

$4$) идеального газа и жидкости
}{
\smallskip\noindent\textbf{Решение}\par
Данная модель лучше всего демонстрирует поведение молекул твердого тела, потому что именно в данном агрегатном состоянии частицы совершают малые колебания вокруг положений равновесия, не перемещаясь по объему тела.

Ответ: $3$
}

\end{document}
